% Propietario preservado y actualizado por el cierre sucesor de sus obligaciones internas.
% Integración por delta: arquitectura cosmológica con frontera, rebote y
% componentes. No altera los teoremas Cartan--Holst ni el selector de G.

\section{Conditions de frontière, régularité du rebond et décomposition des composantes cosmologiques}
\label{chap:cosmologia-frontera-rebote}
\index{cosmologie!réalisation par feuillets}
\index{rebond cosmologique}\index{composante cosmologique descendante}
\index{cobordisme}\index{information!conservation cosmologique}

Ce chapitre distingue deux domaines. La première section démontre les trois
feuilletages de de Sitter. Les sections suivantes définissent les tuples de
données nécessaires à une réalisation cosmologique : sélecteur HMT, carte de
frontière, dynamique de rebond et composante descendante. Chaque conclusion est
restreinte au sous-domaine dans lequel les équations d'évolution et
les conditions de raccordement déclarées sont vérifiées.

\subsection{Correspondance entre les régimes elliptique, parabolique et hyperbolique du TRIT et les cartes AdS–euclidienne–dS}
\label{sec:atlas-ads-euclideo-ds-hmt}

Les trois algèbres quadratiques du TRIT reçoivent une réalisation géométrique
coordonnée. Pour les échelles \(\ell,H^{-1}>0\), nous fixons les cartes
\[
 \mathcal M_-=(\mathbb D, g_{\rm AdS}),
 \qquad
 \mathcal M_0=(\mathscr H_0^{(L)},g_{\rm E}),
 \qquad
 \mathcal M_+=(\mathrm{dS}_4,g_{\rm dS}),
\]
où la section spatiale hyperbolique est représentée dans le disque de Poincaré
par
\[
 g_{\rm AdS}^{\rm sp}
 =\frac{4\ell^2\,|dz|^2}{(1-|z|^2)^2},
 \qquad |z|<1,
\]
le L gnomonique porte la carte euclidienne \(g_{\rm E}\), et la carte positive
adopte l'un des feuilletages de de Sitter construits dans la section
suivante. Le foncteur géométrique conserve l'étiquette de feuillet et les applications de
changement :
\[
 \mathfrak R_{\rm geo}:\mathcal X_{\rm TPK}^{\rm enr}
 \longrightarrow
 \mathcal M_-\cup\mathcal M_0\cup\mathcal M_+.
\]

La complétion EMRH distingue l'intérieur, la couronne et le bord, tandis que la double
projection conserve une même section avant ses évaluations. Sur la
carte négative, la restriction au bord et la représentation d'incidence
forment le carré
\[
\begin{CD}
 \mathcal X_{\rm TPK}^{\rm enr} @>{\mathfrak R_{\rm AdS}}>>
 \mathcal M_-\\
 @V{\mathcal P_{\rm exc}}VV @VV{\operatorname{Res}_{\partial}}V\\
 \mathscr E_{\partial} @>>{\mathfrak R_{\rm conf}}>
 \operatorname{Obs}(\partial\mathcal M_-),
\end{CD}
\]
avec \(\mathfrak R_{\rm conf}\circ\mathcal P_{\rm exc}
=\operatorname{Res}_{\partial}\circ\mathfrak R_{\rm AdS}\). Cette identité
réalise dans HMT la correspondance holographique intérieur--bord : l'état de l'
intérieur et les observables conformes de frontière sont deux représentations de
la même généalogie.

La carte euclidienne centralise la descente locale des lecteurs inertiel et
gravitationnel,
\[
 \mathcal G_L/\mathcal I_L=1,
\]
et constitue la résidence du principe d'équivalence. L'état antérieur
à cette projection conserve l'Ambivalence d'aire--volumique,
\[
 \mathcal L_{\rm ar}/\mathcal L_{\rm vol}
 =\pi_{\rm HMT}/\varphi_{\rm HMT}^{2}.
\]
La carte positive complète l'atlas avec \(\ddot a/a=H^2>0\). Par conséquent, le
même triplet géométrique contient la correspondance holographique en courbure
négative, l'équivalence locale en courbure nulle et l'accélération cosmologique
en courbure positive.

\subsection{\(3\) feuilletages exacts de de Sitter}
\label{sec:tres-foliaciones-de-sitter}

\begin{proposition}[Feuilletages fermé, plat et ouvert]
\label{prop:tres-foliaciones-de-sitter}
Pour \(H>0\), les trois feuilletages FLRW de de Sitter de rayon \(H^{-1}\)
s'expriment, avec \(a_\ast>0\) dans la carte plate, par
\[
\begin{array}{ccc}
\toprule
k&a_k(t)&\text{domaine}\\
\midrule
+1&H^{-1}\cosh(Ht)&t\in\mathbb R\\
0&a_\ast e^{Ht}&t\in\mathbb R\quad\text{(carte plane)}\\
-1&H^{-1}\sinh(Ht)&t>0\quad\text{(carte ouverte)}
\\\bottomrule
\end{array}
\]
et satisfont
\[
\boxed{
\left(\frac{\dot a_k}{a_k}\right)^2+\frac{k}{a_k^2}=H^2,
\qquad
\frac{\ddot a_k}{a_k}=H^2,}
\]
\[
\boxed{
R^{(4)}=12H^2,\qquad
R^{(3)}=\frac{6k}{a_k^2}.}
\]
\end{proposition}

\begin{proof}
Pour \(k=+1\), la première équation utilise
\(\tanh^2+\operatorname{sech}^2=1\). Pour \(k=-1\), elle utilise
\(\coth^2-\operatorname{csch}^2=1\). Dans la carte plate,
\(\dot a/a=H\). Les trois fonctions satisfont
\(\ddot a/a=H^2\), et la formule FLRW de courbure donne
\(R^{(4)}=12H^2\). Chaque tranche spatiale tridimensionnelle possède une courbure
sectionnelle constante \(k/a_k^2\) ; en dimension trois, la contraction du tenseur
de courbure donne donc
\(R^{(3)}=3(3-1)k/a_k^2=6k/a_k^2\).
\end{proof}

Le signe \(k\) type la courbure des sections spatiales, non la
courbure de l'espace-temps complet : les trois cartes ont
\(R^{(4)}>0\). La carte fermée est globale ; les cartes plate et ouverte sont
des patches.

\begin{statusbox}
    Pour une route de retour fermée et le choix spatial déclaré
    \(D=3\), le cocycle HMT d'échelle apporte les facteurs exacts
    \(\varphi^3,\varphi^9,\varphi^{36}\) aux profondeurs
    \(9,27,108\), avec un triplet spectral de mémoire et un écran de
    signature \((3,1)\). La sélection cosmologique consiste à construire, sur
    un état enrichi de TPK, une mesure projective, un facteur d'échelle et une équation d'
    évolution qui entrelacent ces objets avec les feuilletages précédents. Sa
    variable de durée doit être dérivée de la droite d'action du
    \cref{chap:escala-accion} et de sa carte métrologique, non identifiée
    au simple nombre de pas. \(H_0\), les densités cosmologiques et le signe
    de l'accélération ne sont pas employés comme entrées du sélecteur.
\end{statusbox}

Les trois feuilletages constituent le codomaine géométrique exact de
comparaison. Le sélecteur construit sur l'état enrichi détermine la
carte réalisée par chaque régime, compose leurs domaines et conserve les
observables qui distinguent la sélection. L'identification à
\(k=+1,0,-1\) s'effectue au moyen de ce sélecteur et non par la simple coïncidence
de trois signes.

\subsection{Univers à frontière}
\label{sec:universo-con-frontera-hmt}

La complétion
\[
 1000=729+271=729+243+27+1
\]
distingue l'intérieur, les strates de frontière et le sceau. Dans une réalisation
cosmologique, cette stratification oriente le type d'une construction sur
une paire
\[
 (M,\partial M),
\]
munie d'une cotétrade, d'une connexion de Cartan, de matière, de mémoire et de flux. L'
identité cardinale ne produit pas à elle seule une variété, une métrique, une
cotétrade ni des conditions de bord. C'est un guide d'organisation qui n'
acquiert un contenu physique que par une application
\[
 \mathfrak R_{\rm cos}:
 \mathcal X_{\rm TPK}^{\rm enr}
 \dashrightarrow
 \bigl(M,\partial M,e,\omega,\Psi,\mathcal B_{\partial M}\bigr).
\]

La frontière fait partie constitutive du problème variationnel. Il faut
déclarer son caractère spatial, temporel ou nul ; les termes de bord et de
coin ; les données fixées ; et les conditions de compatibilité avec les
équations du mouvement. La mémoire HMT fournit un domaine candidat
pour enregistrer les histoires et les prolongements, mais son identification aux données
géométriques exige l'entrelacement des holonomies du
\cref{crit:puente-tpk-cartan}.

\subsection{Critère de régularité du rebond sous concentration extrême et dominance torsionnelle}
\label{sec:rebote-cosmologico-condicionado}

Une concentration non bornée de densité d'énergie, d'action ou de courbure ne
peut être attribuée à un mécanisme de stockage énergétique dans le secteur
torsionnel. Dans l'action minimale, le courant de spin détermine localement la
contorsion ; la courbure conserve son secteur radiatif. Une correction
répulsive de haute densité requiert un modèle de matière à spin, un
couplage, une équation d'état et un signe effectif déterminés.

Dans un modèle d'Einstein--Cartan avec fluide de spin peut apparaître
\[
 s^2\propto a^{-6}.
\]
Si \(\rho_{\rm m}\) désigne la densité de masse, une équation effective peut
prendre la forme
\[
 H^2
 =
 \frac{8\pi G}{3}
 \bigl(\rho_{\rm m}-\rho_{\rm spin,m}\bigr)
 +\cdots.
\]
Si l'on emploie une densité d'énergie \(\rho_{\rm E}\), la conversion dimensionnelle
oblige à écrire
\[
 H^2
 =
 \frac{8\pi G}{3c^2}
 \bigl(\rho_{\rm E}-\rho_{\rm spin,E}\bigr)
 +\cdots.
\]
Les deux conventions ne sont pas mélangées dans une même équation.

Le candidat au rebond se situe à
\[
 \rho_{\rm m}=\rho_{\rm spin,m}
\]
ou à l'égalité énergétique équivalente, pourvu que la solution soit régulière
et satisfasse les conditions de matière et de frontière déclarées. L'action
Cartan--Holst permet d'étudier ce mécanisme ; HMT doit encore construire l'
application qui sélectionne le courant, le seuil et le registre d'information sans
les consulter rétrospectivement.

\subsection{Cobordisme et composantes cosmologiques descendantes}
\label{sec:cobordismo-universo-hijo}

L'ouverture d'une composante cosmologique descendante se formule au moyen d'un cobordisme
\[
 \mathcal W:
 \Sigma_{\rm in}
 \longrightarrow
 \Sigma_{\rm prog}\sqcup\Sigma_{\rm desc}.
\]
Avant d'affirmer la conservation de l'information, il faut fixer l'objet
conservé. Si \(\mathscr I\) est une charge additive définie par un cocycle,
on peut essayer
\[
 \mathscr I(\Sigma_{\rm in})
 =
 \mathscr I(\Sigma_{\rm prog})
 +\mathscr I(\Sigma_{\rm desc})
 +\mathscr F_{\partial\mathcal W}.
\]
Si \(\mathscr I\) représente de l'information quantique, la structure pertinente
est une isométrie ou une évolution unitaire et le bilan ne se réduit pas, en
général, à une somme scalaire. Il faut calculer les entropies des
sous-systèmes, leurs corrélations et l'information mutuelle.

Une réalisation doit spécifier :
\begin{enumerate}
 \item les conditions de jonction entre les sections initiale, parente et
 descendante ;
 \item les termes de bord et de coin de l'action ;
 \item la causalité, l'hyperbolicité et le problème à valeurs initiales ;
 \item le contrôle des singularités, des courbes temporelles fermées et des conditions
 d'énergie ;
 \item la charge, l'isométrie ou la loi de conservation précise qui représente l'
 information.
\end{enumerate}
Les identités de Bianchi et le bilan de Noether sont des conditions locales
nécessaires ; ils ne démontrent pas à eux seuls le changement de topologie ni la
conservation globale de l'information.

\subsection{Composantes parallèles et mode intercomposantes}
\label{sec:componentes-cosmologicas-paralelas}

Une frontière non connexe et une collection de secteurs possèdent des compositions
distinctes. Dans une réalisation de type topologique,
\[
 \mathcal H_{\Sigma_{\rm parent}\sqcup\Sigma_{\rm child}}
 \simeq
 \mathcal H_{\rm parent}\otimes\mathcal H_{\rm child}
\]
représente la composition des composantes d'une frontière. En revanche,
\[
 \mathcal H_{\rm cos}
 =\bigoplus_a\mathcal H_a
\]
représente des alternatives ou des secteurs de supersélection. Le produit tensoriel
et la somme directe ne se substituent pas l'un à l'autre.

Sur la somme directe, un opérateur global candidat a une forme par blocs :
\[
 \mathbb H_{\rm cos}
 =
 \begin{pmatrix}
 H_1&K_{12}&\cdots\\
 K_{21}&H_2&\cdots\\
 \vdots&\vdots&\ddots
 \end{pmatrix}.
\]
Pour définir un opérateur physique, il faut fixer un domaine dense commun ou des
hypothèses de bornitude relative, une horloge commune ou relationnelle et la condition
\[
 K_{ba}=K_{ab}^{\dagger}.
\]
Les blocs non diagonaux doivent être dérivés de la frontière, du cobordisme ou
d'un autre couplage déclaré. S'ils sont non nuls, les composantes forment des
secteurs couplés d'une théorie globale ; ce ne sont pas des univers indépendants au
sens fort.

Un vecteur propre dont le support s'étend sur plusieurs facteurs de la somme constitue un mode
intercomposantes. Appeler son quantum graviton exige en outre que le secteur
physique réduit possède les hélicités \(\pm2\), deux polarisations transverses,
une énergie positive, une propagation à la vitesse de la lumière, une absence de fantômes et de tachyons,
des identités de Ward et un couplage universel. Le support global ne
satisfait pas à lui seul ces conditions.

\subsection{Certification du rebond torsionnel et obstructions exactes}
\label{sec:condiciones-prueba-cosmologia-hmt}

La réalisation est certifiée par les constructions suivantes :
\begin{enumerate}
 \item des solutions régulières de rebond avec matière et signe dérivés ;
 \item un problème aux limites bien posé ;
 \item un cobordisme admissible avec conditions de jonction ;
 \item un transport d'information défini comme charge, isométrie ou
 évolution unitaire ;
 \item un opérateur global autoadjoint avec des couplages non diagonaux
 sélectionnés avant la comparaison ;
 \item des prédictions qui distinguent la réalisation d'une cosmologie à une
 seule composante.
\end{enumerate}

Le certificat exige la solution régulière, la compatibilité des
conditions de jonction, la conservation de la charge ou de l'isométrie choisie, l'
indépendance par rapport aux cartes arbitraires et la préservation de la causalité
et de l'hyperbolicité. L'énergie globale n'est employée que lorsqu'il existe une
symétrie temporelle ou une définition quasi locale explicite.
